\section{Conditioned matrix estimates}\label{prof:matrices}

The scalar diagnostics now determine the small eigenvalues of the
conditional matrix. The point is to retain the coefficient $1$ in the
lower edge $q^2$, which later yields the coefficient $1$ in $s^{-1}$.

Put
\[
 u=\frac{m}{\sqrt{nq}},\qquad E=u^\perp,\qquad
 e=\frac{P_E\mathbf1}{\sqrt n},\qquad
 c_3=\frac{P_{\operatorname{span}(m,\mathbf1)^\perp}y}
            {\|P_{\operatorname{span}(m,\mathbf1)^\perp}y\|}.
\]
The scalar diagnostics make $c_3$ well defined. The plant component
$e$ is not normalized: $\|e\|^2=1-M^2/q$. Let
$\mathcal P=\operatorname{span}(u,e,c_3)$ and let
$V_{\mathcal P}=(P_{\mathcal P}VP_{\mathcal P})|_{\mathcal P}$.
Operators defined on $\mathcal P$ are extended by zero on its
orthogonal complement. Consider
\[
 \begin{split}
 K&=V^{-1}+bI-J-2q uu^{\mathsf T}/\langle v^2\rangle,\\
 K_o&=(b+b^{-1})I-J+
       [V_{\mathcal P}^{-1}-b^{-1}I_{\mathcal P}]
       -2q uu^{\mathsf T}/b^2,\\
 K_T&=V^{-1}+bI-J-2q uu^{\mathsf T}.
 \end{split}
\]

\begin{proposition}[Conditioned precision bounds]\label{prof:conditioned-precision}
Fix a sufficiently small $\eta>0$, an accuracy $\epsilon>0$, and a
budget $B>0$. On retained scalar data with the required accuracies,
the following estimates hold with the integrated exclusion specified
below. For $\mathsf K\in\{K,K_o,K_T\}$,
\begin{equation}\label{prof:eq10}
 \mathsf K\succeq(1-\epsilon)q^2I,\qquad
 \operatorname{Tr}\mathsf K^{-2}\le Cn/q,\qquad
 \|A_0\mathsf K^{-1/2}\|\le Cq^{-\eta},
 \quad A_0=2I-J.
\end{equation}
For $K_o$ the last bound holds with $\eta=0$. Each precision has at
least $c_A nq^3$ eigenvalues below $C_Aq^2$, and
$\operatorname{Tr}K_o^{-1}=nb+o(\sqrt{n/q})$. For
$\mathsf K\in\{K,K_o\}$, the three projection bounds are
\[
 u^{\mathsf T}\mathsf K^{-1}u\le C_Aq^{-2},\qquad
 e^{\mathsf T}\mathsf K^{-1}e\le C_Aq^{-1},\qquad
 c_3^{\mathsf T}\mathsf K^{-1}c_3\le C_A.
\]

Let $\mathcal S$ be the retained scalar set, let $\mathcal G_{\rm spec}$
be the original spectral event, and let $\mathcal F_y$ be failure of
any estimate above. In the conditional law $\mathsf Q_y$ of
Lemma~\ref{prof:root-count}, the bound is
\[
 \int_{\mathcal S}G(y)\,
       \mathsf Q_y(\mathcal G_{\rm spec}\cap\mathcal F_y)\,dy
 \le n^C e^{-Bk}.
\]
The same bound with the determinant weight follows after increasing
the preliminary budget. The spectral restriction stays inside the
event being estimated; the argument does not condition on its probability.
\end{proposition}

\begin{proof}
\par\smallskip\noindent\textit{Conditional law and Haar concentration.}\enspace
We use Haar concentration in the Frobenius metric on proper rotations:
an $L$-Lipschitz function on $SO(N)$ has tail
$2e^{-cNz^2/L^2}$ about its mean. See
\cite[Sections~5.2--5.3, Proposition~5.13 and Theorem~5.16]{Meckes2019ClassicalGroups}.
A constant-mesh net of size $e^{O(j)}$ then gives an operator bound on
a $j$-dimensional space, using polarization for bilinear forms.
All such tests below condition first on deterministic eigenvalues,
so they also apply to resolvents defined separately on spectral bins.

In $\mathsf Q_y$, the scalar data and the vectors determined by $y$
are fixed. The regression in Lemma~\ref{prof:root-count} gives
\[
 J|_E=W_E+ee^{\mathsf T},\qquad
 p:=P_EJu=p_0+\sqrt{s/\nu}\,w,\qquad
 p_0=\frac{q}{\nu}\frac{P_E(y-d\mathbf1)}{\sqrt{nq}}
                   +\frac{M}{\sqrt q}e,
\]
where $W_E$ is a centered GOE block and $w\sim N(0,I_E/n)$ is
independent of it. The mean $p_0$ is the projected regression mean of
$Wm+\sqrt s\,g$, plus the plant column. We may next condition on the
eigenvalues of $W_E$ and retain its Haar eigenframe. On the intersection
with $\mathcal G_{\rm spec}$, interlacing of
$W_E$, $W_E+ee^{\mathsf T}=J|_E$, and $J$ supplies the required counts
and $\lambda_{\max}(W_E)\le\lambda_{\max}(J)$. The concentration
probabilities are estimated before that intersection is imposed.

\par\smallskip\noindent\textit{Diagonal truncation and the compensating moments.}\enspace
Compare with $D=(b+b^{-1})I-J$. In $K$ and $K_T$, replace the
positive diagonal entry $v_i^{-1}-1$ by
$\min\{v_i^{-1}-1,q^{1-\eta}\}$. This gives a lower precision.
Let $B$ be the difference of either truncated precision, or of $K_o$,
from $D$, and set $B_E=P_EB|_E$ and $f=P_EBu$.

The scalar moments give the quantities needed by the Schur complement.
With
$\mu_B=n^{-1}\operatorname{Tr}B_E$ and
$\nu_B=n^{-1}\operatorname{Tr}B_E^2$,
\[
 \begin{gathered}
 \|B\|=O(q^{1-\eta}),\qquad \mu_B,\nu_B=O(q^2),\qquad
 \mu_B-\nu_B=o(q^2),\\
 u^{\mathsf T}Bu\ge2q^2+o(q^2),\qquad
 \|f\|^2=6q^2+o(q^2),\qquad p_0=f/3+O_A(q^{3/2}),\\
 e^{\mathsf T}f=O_A(q^2),\quad \|B_Ee\|=O(q),\quad
 e^{\mathsf T}B_Ee=O_A(q^2),\quad e^{\mathsf T}p_0=O_A(q^{3/2}).
 \end{gathered}
\]
The error in the vector relation is in Euclidean norm.

For the diagonal precisions, the truncated entry expands as
$m^2+m^4+\cdots$ on $|y|\le q^{1/2-\eta_0}$, with
$0<\eta_0<\eta/2$. The missed bounded powers on the complement are
$O(m^2)$ and have average $O_A(q^4)$ by the tail-square diagnostic.
The bulk mean is therefore $3q^2-q^2=2q^2$ to leading order, equal
to the centered variance $\langle(m^2-q)^2\rangle=2q^2+o(q^2)$.
In the root direction,
\[
 \frac{\langle m^4+m^6\rangle}{q}=3q+7q^2+o(q^2),\qquad
 \frac{2q}{\langle v^2\rangle}=2q+4q^2+o(q^2).
\]
Subtracting the scalar $b^{-1}-1=q+q^2+O(q^3)$ leaves at least
$2q^2+o(q^2)$. Also
\[
 \|P_E\operatorname{diag}(m^2)u\|^2
 =\frac{\langle m^6\rangle}{q}
       -\left(\frac{\langle m^4\rangle}{q}\right)^2
 =6q^2+o(q^2).
\]
This is the leading value of $\|f\|^2$.

The root identity $y=m+m^3/3+O(m^5)$ on the clipped region gives
\[
 \frac{P_Ey}{\sqrt{nq}}
 =\frac13P_E\operatorname{diag}(m^2)u+O_A(q^{3/2}).
\]
The tail contributes $O_A(q^4/q)$ to its squared error. In the formula
for $p_0$, the extra $e$ coefficient is
$(M-(q/\nu)d)/\sqrt q=O_A(q^{3/2})$, giving $p_0=f/3+O_A(q^{3/2})$.
The bounds involving $e$ follow similarly from
$\langle m^3\rangle=3qM+O_A(q^3)$.

For $K_o$, put
$L_B=(P_{\mathcal P}(\operatorname{diag}(m^2)-qI)P_{\mathcal P})|_{\mathcal P}$.
Its norm is $O_A(q)$: in the cubic direction the numerator with an
extra factor $m^2$ is $O(q^4)$ and the normalization is of order
$q^3$. Since $V_{\mathcal P}=bI_{\mathcal P}-L_B$,
\[
 V_{\mathcal P}^{-1}
 =b^{-1}I_{\mathcal P}+b^{-2}L_B+b^{-3}L_B^2+O_A(q^3).
\]
Here $(L_B)_{uu}=2q-8q^2+o(q^2)$ and
$(L_B^2)_{uu}=10q^2+o(q^2)$. Projection onto $\mathcal P$ in the
second identity loses only $O_A(q^3)$, by the preceding approximation
of the cubic column. Substitution gives the same root coefficient and
column estimates; the $e$ entries again use the third moment.

Finally, the diagonal restored in $K$ and $K_T$ is positive and is
supported where $v_i^{-1}-1>q^{1-\eta}$. The tail-square bound limits
this support to $O(nq^{3+\eta'})$ sites for some fixed $\eta'>0$.
Thus restoring it is a positive perturbation of rank $o(nq^3)$.

\par\smallskip\noindent\textit{Resolvents away from the edge.}\enspace
Augment by one scalar $\beta$ for the negative spike $-ee^{\mathsf T}$,
so the augmented form contains
$\beta^2-2\beta e^{\mathsf T}x_E$. Put
$\Delta_{\mathrm{sp}}=q^{8/5}$ and define
\[
 N=\operatorname{ran}\mathbf1_{\{\,2-W_E\le\Delta_{\mathrm{sp}}\,\}},
 \qquad H=E\ominus N.
\]
Then $\dim N=O(n\Delta_{\mathrm{sp}}^{3/2})$. The high-block resolvent
for the unperturbed comparison, also after subtracting
$(1-\epsilon)q^2$, is
\(R=(2+O(q^2)-W_E)^{-1}_H\).

Its traces divided by \(n\) for the first and second powers are \(1+o(1),\,\asymp \Delta_{\mathrm{sp}}^{-1/2}\). On the conditional Haar orientation with the spectra satisfying the interlacings, except at arbitrarily high rate,
\[
 \|R^{1/2}B_E R^{1/2}\|=o(1).
\]
Indeed dyadic blocks at depths $t_1,t_2\ge\Delta_{\mathrm{sp}}$ give
centered bilinear deviations at most
$Cq^{1-\eta}\max(t_1,t_2)^{3/4}$ by Haar concentration and nets.
After division by $\sqrt{t_1t_2}$, the dyadic sum is bounded by a
logarithmic factor times
$q^{1-\eta}\Delta_{\mathrm{sp}}^{-1/2}
=q^{1/5-\eta}$. The mean costs
$O(q^2/\Delta_{\mathrm{sp}})=O(q^{2/5})$. Both vanish for small
$\eta$; the sum includes a constant-width bulk bin.

\par\smallskip\noindent\textit{Uniform tests on the remaining coordinates.}\enspace
For remaining variables \(\gamma,x_N,\beta\), normalized by \(\gamma^2+\|x_N\|^2+\beta^2/\sqrt{\Delta_{\mathrm{sp}}}\le1\), write \(x=x_N+R(p\gamma+e\beta)\). The needed simultaneous Haar estimates are
\begin{equation}\label{prof:eq11}
\begin{split}
 \|x\|&=O_A(1),\quad x^TB_E x=\mu_B\|x\|^2+o(q^2),\quad
 f^T x= (2q^2+o(q^2))\gamma+o(q^2),\\
 \|R^{1/2}(B_E x+f\gamma)\|^2&=\nu_B\|x\|^2+6q^2\gamma^2+o(q^2).
 \end{split}
\end{equation}
Here little-oh can mean any specified fixed leading error.

For a fixed symmetric weight $T$, normalized Gaussian quadratic tails
and polarization give error
\[
 C_B\frac{\sqrt k}{n}\sqrt{\operatorname{Tr}T^2}
       +C_B\frac{k}{n}\|T\|
\]
for a Haar quadratic test on fixed orthonormal axes, at failure
$e^{-Bk}$. The same bound holds for an independent Gaussian column
with covariance $I_E/n$. A subspace test instead pays its dimension
through a net. Since $q\asymp_A r$, fixed changes of the constants
allow $nq^3$ in place of $k$.

Fix, before sampling $W_E$ and $w$, a fixed-dimensional space $S$
containing $e,c_3,p_0,f$. On $S$, compression of $R$ differs from its
trace mean by
$O_A(q^{3/2}\Delta_{\mathrm{sp}}^{-1/4}+q^3/\Delta_{\mathrm{sp}})$;
compression of $R^2$ has relatively vanishing error. Test $w$
separately, with zero mixed means, so $S$ stays deterministic. Also
$\|P_N|_S\|=O(\Delta_{\mathrm{sp}}^{3/4})$.
These estimates give $\|x\|=O_A(1)$ and
\[
 f^{\mathsf T}Rp
 =\frac{\operatorname{Tr}R}{n}f^{\mathsf T}p_0+o(q^2)
 =2q^2+o(q^2).
\]
The terms $f^{\mathsf T}x_N$ and $\beta f^{\mathsf T}Re$ are $o(q^2)$
on the normalized variable ball. This proves the $f^{\mathsf T}x$
line of \eqref{prof:eq11}. The same tests give
\[
 P_Sx=\frac{\operatorname{Tr}R}{n}\,\beta e+O_A(q).
\]

The fixed-axis estimates have identified the projection of the
optimizing vector. We next expose the low frame and the two resolved
columns. Rotations preserving these data control the remaining
components and hence the final weighted quadratic test.

Now expose $N,Rp,Re$ and put
$Y=S+N+\operatorname{span}(Rp,Re)$. Before this exposure, the part of
$Y$ beyond $S$ is invariant under rotations fixing $S$, hence is a
uniform subspace of $S^\perp$. Its compressions of $B_E^j$, $j=1,2$,
have deviations from their scalar means of size
$O_A(q^{j(1-\eta)}\Delta_{\mathrm{sp}}^{3/4})$; mixed compressions
with $S$ have the same bound and mean zero. Together with the estimates
on $B_Ee$, $e^{\mathsf T}B_Ee$, and $P_Sx$, this gives
\[
 x^{\mathsf T}B_Ex=\mu_B\|x\|^2+o(q^2),\qquad
 \|B_Ex\|^2=\nu_B\|x\|^2+o(q^2),
\]
as well as
$\|P_SB_Ex\|=o(q)$ and
$\|P_{Y\cap S^\perp}B_Ex\|=O_A(q^2)$.

After the exposure, only rotations fixing $Y$ pointwise remain.
Split $R=R_{\rm near}+R_{\rm far}$ at a small fixed depth $\zeta$.
The conditional mean of the far restriction on $Y^\perp$ is its
trace on that complement divided by its dimension. Since
$\dim Y=O(n\Delta_{\mathrm{sp}}^{3/2})$,
$\|R_{\rm far}\|\le C/\zeta$, and
$n^{-1}\Tr R_{\rm near}=O(\sqrt\zeta)$, this scalar is
\[
 1+O(\sqrt\zeta)+O(\Delta_{\mathrm{sp}}^{3/2}/\zeta)+o(1).
\]
Residual Haar concentration and the low-dimensional nets make this
the quadratic coefficient on the needed vectors in $Y^\perp$.
The component of $B_Ex$ in $Y$ contributes $o(q^2)$ to the far
quadratic and mixed terms, by its two projection bounds and the fixed
far norm. The mixed term with $Rf$ has zero leading conditional mean;
$f^{\mathsf T}Rf$ is already covered by the fixed-axis test.

For the near part, $\|R_{\rm near}\|\le C/\Delta_{\mathrm{sp}}$ and
its compression to $S$ is small by the fixed-axis tests. Its action
on the $O_A(q^2)$ component in $Y\cap S^\perp$ is negligible at
scale $q^2$. On the remaining vectors, the conditional square-root
norm is bounded by a constant times
\[
 \sqrt{\Tr R_{\rm near}/n}
 +\sqrt{\dim Y/n}\,\|R_{\rm near}^{1/2}\|
 =O(\zeta^{1/4}+\Delta_{\mathrm{sp}}^{1/4}).
\]
This is the Haar Lipschitz bound plus its net cost. First choose
$\zeta$ small, then the upper $q$ scale. These estimates prove the
last line of \eqref{prof:eq11}.

All these failures were estimated under the Gaussian/Haar conditional
law before imposing $\mathcal G_{\rm spec}$. On its intersection,
interlacing supplies the spectra of $W_E$ used in the trace and count
bounds. Integration in $y$ gives the stated exclusion.

\par\smallskip\noindent\textit{Schur complements, counts, and projections.}\enspace
Substitution of the simultaneous tests into the high-block Schur
complements gives the perturbed-minus-unperturbed difference
\[
 \gamma^2 u^TB u+2\gamma f^Tx+x^T B_E x-\|R^{1/2}(B_E x+f\gamma)\|^2.
\]
By \eqref{prof:eq11} and $u^{\mathsf T}Bu\ge2q^2+o(q^2)$, the displayed
difference is bounded below on the normalized variable ball by
\[
 (2+4-6)q^2\gamma^2+(\mu_B-\nu_B)\|x\|^2-o(q^2)=-o(q^2).
\]
This is the one-sided root-direction cancellation and the bulk cancellation.

To see the remaining positive margin, let
$x_{\rm phys}=\gamma u+x_E$. The shifted augmented base form is exactly
\[
 x_{\rm phys}^{\mathsf T}
     \bigl(D-(1-\epsilon)q^2I\bigr)x_{\rm phys}
       +(\beta-e^{\mathsf T}x_E)^2.
\]
The original spectral bound makes its first term at least a fixed
multiple of $\epsilon q^2\|x_{\rm phys}\|^2$. At the high optimum,
$x_H=R(p\gamma+e\beta)$, and
\[
 \|x_H\|^2\ge c\beta^2\|Re\|^2-C\gamma^2\|Rp\|^2,
 \qquad \|Re\|^2\asymp\Delta_{\mathrm{sp}}^{-1/2},\quad
 \|Rp\|=o(1).
\]
After absorption, the base complement therefore controls
$c\epsilon q^2(\gamma^2+\|x_N\|^2+
\beta^2/\sqrt{\Delta_{\mathrm{sp}}})$. The $o(q^2)$ difference
preserves this margin and proves the lower precision bound.

The one-sided estimate proves positivity. The trace and shell-density
arguments also need enough eigenvalues near the critical edge. Removing
the fixed-rank directions makes the bulk cancellation two-sided, so the
same calculation can now be repeated at dyadic thresholds.

After removing the fixed-rank $u,e$ interactions, the unaugmented
$N$-Schur complement for a truncated precision differs from the base
one by $o(q^2)$ on both sides, because $\mu_B-\nu_B=o(q^2)$.
At a shifted threshold $a\le\Delta_{\mathrm{sp}}/2$ the high block
remains positive. Schur inertia therefore identifies the number of
full eigenvalues below $a$ with the negative-eigenvalue count of the
low complement, up to the removed fixed ranks. The two-sided comparison
bounds this count by the $W_E$ edge counts at thresholds displaced by
$o(q^2)$.

The same estimates hold on a finite dyadic grid
$Cq^2\le a\le\Delta_{\mathrm{sp}}/2$: the shift is at most half the
high gap, so the high resolvents remain comparable and their normalized
traces change by $O(\sqrt{\Delta_{\mathrm{sp}}})$. This gives the
upper count $Cna^{3/2}$ and at least $cnq^3$ eigenvalues below
$Cq^2$ for a sufficiently large fixed $C$. Above a fixed multiple of
$\Delta_{\mathrm{sp}}$, the weighted dyadic form bound gives the same
upper count. Restoring the positive discarded diagonal preserves upper
counts and removes at most its rank $O(nq^{3+\eta'})=o(nq^3)$ from
the lower spread. Dyadic integration of these counts gives all
trace-square bounds.

For $K_o$, the correction to $D$ has fixed rank. At
$b+b^{-1}$, the semicircle transform is exactly $b$;
\eqref{prof:eq1} and the fixed-rank trace error $O(q^{-2})
=o(\sqrt{n/q})$ give its stated inverse trace. For a truncated
precision, write $A_0$ as that precision minus an
$O(q^{1-\eta})$ term and use the lower bound at $q^2$ to obtain
$\|A_0\mathsf K^{-1/2}\|\le Cq^{-\eta}$. Restoring the positive
diagonal preserves this bound by inverse order. The fixed-rank
structure gives the stated $\eta=0$ bound for $K_o$.

It remains to control variances in the distinguished directions. For
projection bounds at \(e,c_3\) first use \(D\). At \(R_1=(b+b^{-1}-W_E)^{-1}\), deterministic fixed-axis compressions have means \(1+O(q)\) and errors \(O_A(q)\), \(\|R_1 e\|^2\asymp_A 1/q\) (lower via modes at sufficiently enlarged width to get any desired rate). Thus \(1-e^T R_1 e\gtrsim_A q\) by the principal-block top eigenvalue bound at \(2+q^2/2\) and shifting. Also \(e^T R_1 p=O_A(q^{3/2})\), \(c_3^TR_1 p=O_A(q)\). Sherman-Morrison for the spike followed by inclusion of \(u\) (whose inverse variance bound after inclusion is the global operator bound) gives the orders stated.

For a truncated precision or $K_o$, the reduced Schur form controls
the base reduced form up to a factor two. View a projection variance
as the square of the dual norm of its physical source in the augmented
positive form. The high precision stays comparable to $R^{-1}$, and
its minimizing vector changes by
$-(R^{-1}+B|_H)^{-1}P_H(B_Ex+f\gamma)$.
For a source at $e$ or $c_3$, \eqref{prof:eq11} and its bounded
$R^{1/2}$ norm bound the change in testing this optimum by
\[
 C_Aq\bigl(|\gamma|+\|x_N\|+
                    \Delta_{\mathrm{sp}}^{-1/4}|\beta|\bigr).
\]
The reduced energy controls $q^2$ times the square of this
parenthesis, so the additional squared dual norm is $O_A(1)$.
The remaining high square also costs $O_A(1)$. Adding these costs to
the base variances preserves the orders $q^{-1}$ for $e$ and $1$ for
$c_3$; the global lower precision bound supplies the $q^{-2}$ order
for $u$. Inverse order again permits restoration of the diagonal.

\end{proof}

\begin{corollary}[Uniqueness of the root on the retained event]
\label{prof:unique-root}
On the retained field event every root is regular and the root is unique.
The excluded root-count mass, and hence the exceptional field probability
under the spectral restriction, have the chosen exponential rate.
\end{corollary}

\begin{proof}
First remove the null set of critical target fields from
Lemma~\ref{prof:root-count}. For the remaining fields every root is
regular, and the probability of an excluded root is bounded by its
expected count. On the complementary field event,
$V^{1/2}K_TV^{1/2}
=I+V^{1/2}(bI-J-2mm^{\mathsf T}/n)V^{1/2}$ is positive.
The Jacobian $D\Phi_J$ is similar to this matrix and thus has positive
determinant at every root.

Properness and the identity homotopy outside a sufficiently large ball
give degree one. Hence the signed root count is one. With all signs
positive there is exactly one root. More generally, for a regular
target the numbers of positive and negative roots obey
$N_+-N_-=1$, so the excess multiplicity is
$N_++N_--1=2N_-$. This identity also bounds multiplicity errors by
the excluded negative-root count.

On the retained event the unique root is a measurable function of
$(J,h)$; extend it by zero elsewhere when a globally defined profile
is needed. Conditional on $(J,h)$ in the planted experiment, the
planted spin has law $\mu_h$. Thus an equivariant test involving this
profile may replace the plant by a posterior replica, with the
excluded event charged separately. This is a conditional posterior
identity under planted disorder; the later likelihood comparison
performs the transfer to ordinary disorder.
\end{proof}
